% Page 092

\seite{92}

\margmark{17. Mai      A}%
m Montag, den 13.\ d. Mts., wurde der Unterricht\ob
von 10--12 Uhr durch Herrn Schulrat Mahrbach besucht.\ob
Die Pfingstferien dauern vom 18.\ bis ein=\ob
schließlich 27.\ Mai.

\margmark{28. Mai      D}%
ie herrlichen Maientage und der schöne\ob
Mai Blütenschmuck vergebenen Pfingstfest\ob
sind bereitwillig Entsetzlich hat wieder heftig be=\ob
gonnen.

\margmark{15. Juni     T}%
rotz der Vorherzeit im Frühling und Vorsommer=\ob
blick, die Frühjahrzeit wohl hinter den Erwartungen zurück=\ob
stellt, und die Fronleichnamsprozession ist gefeiert worden.

\margmark{1. Juli      I}%
n Reinigungs- Raumschluß\unsicher der Pfarrei Bettingen\ob
sollte am 30.\ Juni in Bettingen sein rechte Räumungs=\ob
fest begehen. Am 25.\ Juni soll gleichzeitig die Nachricht,\ob
ein, daß die am 13.\ Juni von der Festungsbereich\ob
alle Weil Einzubringen zurückgezogen sei, die Feitstunden\unsicher\ob
Feier konnte demzufolge nicht gefallen werden.\ob
Mit heutigem Tag begannen die Ferien, sie\ob
dauern 14 Tage.

\margmark{20. Juli     W}%
egen mangelhafter Restbestände sind in die\ob
Ferien am 3.\ und 4.\ Juli unterbrochen, der Unter=\ob
richt daher erst wieder am 18.\ Fortgenommen.\ob
Zwei schwere Gewitter zogen heute nacheinander\ob
zwischen 8 und 8 Uhr abends durch den Ort. Der erste\ob
Blitz schlag die Licht zuerst in das Haus der Familie\ob
Müller ohne zum Glück zu zünden. Ein zweiter\ob
Blitzschlag, der bald darüber folgte, halber zunächst\ob
wohl die Dachstühle\unsicher erhalten, da aber viel Heu, Mehl, Lehm,\ob
sowie mit zwei Fuhren Getreide, Stroh und\ob
Peter Goebel. Die Feuer wurden sofort, bei wechselnd\ob
der Festmanndienst das Jahr zehnt Geschädigten. Eine\ob
Hilfe und Retter und des Baurhof,\unsicher{} bei Konto\ob
da er alle Tat\unsicher der Einschlag sich den Wagenträger\ob
an der Bekanntmachung 6 Stuhl. In den zweiten\ob
Gewitterschlag zog das Licht ableitend ein, und zumal\ob
in die von den Hauß des Joh.\ Weimarg\unsicher Gegenteil\unsicher\ob
Nagel an die andere Seite des Hauses warnen.

\margmark{5. August    D}%
ie Getreideernte ist voll im Gange. Viel Korn\ob
in der Jahrhundertgenannten\unsicher eingetrokknetem Regen. Wird\ob
nur ungefähr vom Landwirt gewählt\unsicher.
\pend
